\chapter{Cloud computing, container e Docker}
\label{cap:container}

\begin{inbreve}
Il cloud computing mette a disposizione risorse virtualizzate su richiesta. Per sfruttarlo il software deve poter
girare in ambienti isolati e portabili. Il capitolo confronta le macchine virtuali con i container e presenta
Docker (immagini, container, registry, comandi). Mostra poi due esempi completi: un'applicazione con un solo
container e un'applicazione a più servizi gestita con Docker Compose. I modelli di servizio e l'architettura del
software cloud sono nel capitolo~\ref{cap:cloud}.
\end{inbreve}

\section{Cloud computing}
\label{sec:container-cloud}

Hardware potente e reti veloci hanno portato allo sviluppo del \termine{cloud computing}: risorse virtualizzate
(calcolo, memoria di massa, piattaforma) accessibili su richiesta. I vantaggi principali sono quattro.

\begin{itemize}
  \item \textbf{Scalability}: mantenere le prestazioni al crescere del carico.
  \item \textbf{Elasticity}: adattare la configurazione dei server alle richieste che cambiano.
  \item \textbf{Resilience}: mantenere il servizio in caso di guasto dei server.
  \item \textbf{Cost}: riduzione delle spese di capitale (\ingl{CapEx}) e pagamento in base all'uso
    (\ingl{pay-per-use}).
\end{itemize}

\section{Dalle macchine virtuali ai container}
\label{sec:container-vm}

Nella \termine{virtualizzazione dei server} molte macchine virtuali (VM), ciascuna con il proprio sistema operativo,
girano sullo stesso server fisico. Un \termine{hypervisor} di ``tipo 1'' è caricato direttamente sull'hardware e
gestisce le VM.

I \termine{container} sfruttano invece la capacità del kernel del sistema operativo di permettere più istanze isolate
dello spazio utente (\figrif{fig:container-vm}). Rispetto alle VM:
\begin{itemize}
  \item[$+$] sono più leggeri: richiedono meno risorse (da Gb a Mb);
  \item[$+$] partono più in fretta (da minuti a secondi);
  \item[$-$] condividono lo stesso sistema operativo.
\end{itemize}

\begin{figure}[ht]
\centering
\begin{tikzpicture}[r/.style={draw, minimum width=5.4cm, minimum height=0.7cm, font=\small, align=center},
  a/.style={draw, minimum width=2.5cm, minimum height=1.3cm, font=\small, align=center, fill=grigiochiaro}]
\node[font=\small\bfseries\sffamily] at (0,4.6) {Macchine virtuali};
\node[r, fill=grigio!30] (h1) at (0,0) {Server hardware};
\node[r] (hy) at (0,0.8) {Hypervisor};
\node[a] at (-1.45,2.2) {App A\\Bins/Libs\\Guest OS};
\node[a] at (1.45,2.2) {App B\\Bins/Libs\\Guest OS};
\node[font=\small\bfseries\sffamily] at (8,4.6) {Container};
\node[r, fill=grigio!30] (h2) at (8,0) {Server hardware};
\node[r] (os) at (8,0.8) {Host OS};
\node[r] (cm) at (8,1.6) {Container manager};
\node[a, minimum height=1cm] at (6.55,2.55) {App A\\Bins/Libs};
\node[a, minimum height=1cm] at (9.45,2.55) {App B\\Bins/Libs};
\end{tikzpicture}
\caption{Pile a confronto. A sinistra ogni VM contiene un sistema operativo ospite completo (\emph{Guest OS}) sopra
l'hypervisor. A destra i container contengono solo l'applicazione e le sue librerie e condividono il sistema operativo
(\emph{Host OS}) sotto il container manager.}
\label{fig:container-vm}
\end{figure}

\section{Docker}
\label{sec:container-docker}

\begin{definizione}[Docker]
\termine{Docker} è una piattaforma che permette di eseguire applicazioni in un ambiente isolato e di sviluppare ed
eseguire applicazioni portabili sfruttando i container.
\end{definizione}

Docker usa la virtualizzazione basata sui container per eseguire più istanze ospiti isolate sullo stesso sistema
operativo. La piattaforma è composta da:
\begin{itemize}
  \item \textbf{Docker Engine}, per creare ed eseguire i container;
  \item \textbf{Docker Hub}, per distribuirli.
\end{itemize}

I concetti principali sono tre:
\begin{itemize}
  \item i componenti software sono impacchettati in \termine{immagini} (\ingl{image}), che sono modelli di sola
    lettura usati per creare ed eseguire i \termine{container};
  \item i container girano sopra il Docker Engine e il sistema operativo del server;
  \item si possono montare \termine{volumi} esterni (\ingl{volume}) per rendere persistenti i dati.
\end{itemize}

\section{Immagini, registry e comandi}
\label{sec:container-immagini}

Le immagini sono modelli di sola lettura. Sono conservate in un \termine{registry} Docker, privato o pubblico
(per esempio Docker Hub).
\begin{itemize}
  \item Il registry è strutturato in \emph{repository}.
  \item Ogni repository contiene un insieme di immagini, per le diverse versioni di un software.
  \item Un'immagine è identificata da una coppia \texttt{repository:tag}.
\end{itemize}

Le immagini sono strutturate in \emph{strati} (\ingl{layer}) di sola lettura (\figrif{fig:container-strati}). Ogni
strato è a sua volta un'immagine. Lo strato più basso si chiama \termine{base image}. Un container in esecuzione scrive
in un nuovo strato in cima, e le modifiche possono essere salvate (\ingl{commit}) in una nuova immagine.

\begin{figure}[ht]
\centering
\begin{tikzpicture}[l/.style={draw, minimum width=4.2cm, minimum height=0.65cm, font=\small}]
\node[l, fill=grigiochiaro] (a) at (0,0) {Debian (base image)};
\node[l, fill=grigiochiaro] (b) at (0,0.65) {Java};
\node[l, fill=grigiochiaro] (c) at (0,1.3) {Dropwizard};
\node[l, fill=white, thick] (d) at (0,1.95) {Container (strato scrivibile)};
\node[font=\footnotesize, anchor=west, align=left] at (2.4,0.65) {strati dell'immagine:\\sola lettura};
\draw[grigio] (2.15,-0.3) -- (2.3,-0.3) -- (2.3,1.6) -- (2.15,1.6);
\end{tikzpicture}
\caption{Struttura a strati di una immagine e di un container. Gli strati dell'immagine (Debian, Java, Dropwizard)
sono di sola lettura. Il container aggiunge in cima uno strato in cui può scrivere.}
\label{fig:container-strati}
\end{figure}

I comandi principali (\figrif{fig:container-comandi}) sono:
\begin{description}[leftmargin=0pt, style=unboxed]
  \item[build] crea un'immagine a partire da un \emph{Dockerfile};
  \item[run] crea ed esegue un container da un'immagine;
  \item[commit] salva lo stato di un container in una nuova immagine;
  \item[push] invia un'immagine a un registry;
  \item[pull] scarica un'immagine da un registry.
\end{description}

\begin{figure}[ht]
\centering
\begin{tikzpicture}[b/.style={draw, minimum width=2.2cm, minimum height=0.8cm, font=\small}]
\node[b] (d) at (0,0) {Dockerfile};
\node[b, fill=grigiochiaro] (i) at (4.4,0) {image};
\node[b] (r) at (8.8,0) {registry};
\node[b] (c) at (4.4,-2.4) {container};
\draw[freccia] (d) -- node[above, font=\footnotesize] {build} (i);
\draw[freccia] ([yshift=3pt]i.east) -- node[above, font=\footnotesize] {push} ([yshift=3pt]r.west);
\draw[freccia] ([yshift=-3pt]r.west) -- node[below, font=\footnotesize] {pull} ([yshift=-3pt]i.east);
\draw[freccia] ([xshift=-4pt]i.south) -- node[left, font=\footnotesize] {commit} ([xshift=-4pt]c.north);
\draw[freccia] ([xshift=4pt]c.north) -- node[right, font=\footnotesize] {run} ([xshift=4pt]i.south);
\end{tikzpicture}
\caption{Ciclo di vita delle immagini Docker. Il Dockerfile produce un'immagine (\texttt{build}); l'immagine si
invia al registry e si scarica da esso (\texttt{push}, \texttt{pull}); da un'immagine si avvia un container
(\texttt{run}) e lo stato di un container diventa una nuova immagine (\texttt{commit}).}
\label{fig:container-comandi}
\end{figure}

\section{Esempio: un'applicazione con un container}
\label{sec:container-hello}

L'esempio è un'applicazione PHP che mostra ``Hello, World''. Ci vogliono due file. Il primo è \texttt{index.php}:

{\small
\begin{verbatim}
<?php
echo "Hello, World"
\end{verbatim}
}

Il secondo è il \texttt{Dockerfile}, che descrive come costruire l'immagine. Si parte da un'immagine PHP trovata su
Docker Hub:

{\small
\begin{verbatim}
FROM php:7.0-apache
COPY src/ /var/www/html
EXPOSE 80
\end{verbatim}
}

Poi si danno due comandi:

{\small
\begin{verbatim}
docker build -t hello-world .
docker run -p 80:80 hello-world
\end{verbatim}
}

\subsubsection{Che cosa fanno istruzioni e comandi\aggiunta}
\texttt{FROM} sceglie l'immagine di partenza, \texttt{COPY} copia la cartella \texttt{src/} nella cartella servita dal
web server, \texttt{EXPOSE} dichiara la porta 80 del container. Il comando \texttt{docker build} costruisce
l'immagine a partire dal Dockerfile della cartella corrente (\texttt{.}) e le dà il nome \texttt{hello-world}
(\texttt{-t}). Il comando \texttt{docker run} avvia un container e collega la porta 80 del computer ospite alla porta
80 del container (\texttt{-p}).

Due regole pratiche:
\begin{itemize}
  \item si può usare un \emph{volume} per condividere una cartella tra il computer ospite e il container;
  \item si usa \emph{un solo processo per container}.
\end{itemize}

\section{Applicazioni a più servizi: Docker Compose}
\label{sec:container-compose}

Come si gestisce con Docker un'applicazione fatta di più (micro)servizi? Docker Compose permette di descriverla in un
solo file e di avviarla con un solo comando. L'esempio ha due servizi: un \emph{product service}, scritto in Python,
che restituisce un elenco di prodotti, e un \emph{website}, scritto in PHP, che interroga il servizio e mostra
l'elenco.

\subsection{Il servizio dei prodotti}

\texttt{product/api.py}:

{\small
\begin{verbatim}
from flask import Flask
from flask_restful import Resource, Api

app = Flask(__name__)
api = Api(app)

class Product(Resource):
    def get(self):
        return {
            'product': ['Ice cream',
                        'Chocolate',
                        'Fruit']
        }

api.add_resource(Product, '/')

if __name__ == '__main__':
    app.run(host='0.0.0.0', port=80, debug=True)
\end{verbatim}
}

\texttt{product/Dockerfile}:

{\small
\begin{verbatim}
FROM python:3-onbuild
COPY . /usr/src/app
CMD ["python", "api.py"]
\end{verbatim}
}

\subsection{Il file di composizione}

Con Compose non servono i comandi \texttt{docker build} e \texttt{docker run} per ogni servizio. Basta il file
\texttt{docker-compose.yml}:

{\small
\begin{verbatim}
version: '3'

services:
  product-service:
    build: ./product
    volumes:
      - ./product:/usr/src/app
    ports:
      - 5001:80

  website:
    image: php:apache
    volumes:
      - ./website:/var/www/html
    ports:
      - 5000:80
    depends_on:
      - product-service
\end{verbatim}
}

\subsubsection{Che cosa dice il file\aggiunta}
Il file dice, per ogni servizio, come ottenere l'immagine (\texttt{build} da una cartella, oppure \texttt{image} da
Docker Hub), quali volumi montare, quali porte collegare (\texttt{ospite:container}) e, per il sito, da quale
servizio dipende (\texttt{depends\_on}). Il comando che avvia tutto è:

{\small
\begin{verbatim}
docker-compose up
\end{verbatim}
}

\subsection{Il sito}

\texttt{website/index.php} legge dal servizio la lista dei prodotti e la stampa. Compose crea una \emph{rete
virtuale}: per questo il sito può usare direttamente il \emph{nome del servizio} (\texttt{product-service}) come
indirizzo.

{\small
\begin{verbatim}
<html>
  <head>
    <title>My Shop</title>
  </head>

  <body>
    <h1>Welcome to my shop</h1>
    <ul>
      <?php
        $json = file_get_contents('http://product-service');
        $obj = json_decode($json);

        $products = $obj->product;
        foreach ($products as $product) {
          echo "<li>$product</li>";
        }
      ?>
    </ul>
  </body>
</html>
\end{verbatim}
}

Aprendo \texttt{localhost:5001} si vede la risposta JSON del servizio. Aprendo \texttt{localhost:5000} si vede la pagina
del sito, che elenca i prodotti.

\begin{puntichiave}
\begin{itemize}
  \item Cloud computing: risorse virtualizzate su richiesta. Vantaggi: scalability, elasticity, resilience, cost.
  \item Container rispetto alle VM: più leggeri (da Gb a Mb) e più veloci ad avviarsi (da minuti a secondi), ma
    condividono il sistema operativo.
  \item Docker: le immagini (sola lettura, a strati) sono i modelli da cui si creano i container. I volumi rendono
    i dati persistenti.
  \item Un'immagine ha un nome \texttt{repository:tag} ed è conservata in un registry (per esempio Docker Hub).
  \item Comandi: \texttt{build}, \texttt{run}, \texttt{commit}, \texttt{push}, \texttt{pull}.
  \item Docker Compose descrive in un file un'applicazione a più servizi e la avvia con \texttt{docker-compose up}.
    I servizi si raggiungono per nome.
\end{itemize}
\end{puntichiave}

\begin{domande}
\begin{enumerate}
  \item Quali sono i quattro vantaggi principali del cloud computing?
  \item Che differenza c'è tra una VM e un container? Quali vantaggi e svantaggi hanno i container?
  \item Che cosa sono un'immagine, un container, un volume, un registry?
  \item Che cos'è un'immagine ``a strati'' e dove scrive un container in esecuzione?
  \item A che cosa servono i comandi \texttt{build}, \texttt{run}, \texttt{commit}, \texttt{push}, \texttt{pull}?
  \item A che cosa serve Docker Compose e come fa il sito a raggiungere il servizio dei prodotti?
\end{enumerate}
\end{domande}
